\subsection{自由膨胀量怎样转化为约束力}
先区分两种应变。温度变化引起的自由热应变是 $\alpha\Delta T$；外力引起的轴向弹性应变是 $\sigma/Y$，其中 $Y$ 为杨氏模量，规定拉应力为正。在小形变、线弹性和温度均匀的近似下，总应变为
\[
 \varepsilon=\alpha\Delta T+\frac{\sigma}{Y}.
\]
固定两端不是令每一项都为零，而是要求总伸长为零。于是均匀杆升温后有 $\sigma=-Y\alpha\Delta T$，负号表示受压。前面所说的热应力量值 $Y\alpha\Delta T$，应与这里的有符号应力区分。
若杆可以自由移动，则约束力为零，伸长才等于自由膨胀量；若已经屈服、脱离支座或发生屈曲，就要重新建模。

\begin{origexamplebox}{两种串联材料受热后的共同轴力}\label{ch14:ex:teaching-composite}
两根细杆首尾牢固连接，长度为 $L_1,L_2$，截面积同为 $A$，杨氏模量为 $Y_1,Y_2$，线膨胀系数为 $\alpha_1,\alpha_2$。初温时恰好无应力地安装在两个固定支座间。现均匀升温 $\Delta T>0$，求轴力大小和两杆各自的长度变化。假定保持直线，不发生屈曲；两个外端支座固定，两杆在连接处无相对滑移，但连接点可以整体移动。
\end{origexamplebox}
\begin{xesolution}
取拉力 $F$ 为正。连接处没有额外轴向外力，静力平衡要求两杆承受同一轴力，而不是承受相同伸长。分别有
\[
 \Delta L_i=\alpha_iL_i\Delta T+\frac{FL_i}{AY_i},\qquad i=1,2.
\]
固定的是两杆总长，因此 $\Delta L_1+\Delta L_2=0$，得到
\[
 F=-A\Delta T\frac{\alpha_1L_1+\alpha_2L_2}{L_1/Y_1+L_2/Y_2}.
\]
若两膨胀系数均为正，$F<0$，两杆受压，压力大小为 $|F|$。将同一个 $F$ 代回第一式，并将分母同乘 $Y_1Y_2$，得
\[
 \Delta L_1=\frac{L_1L_2(\alpha_1Y_1-\alpha_2Y_2)}{L_1Y_2+L_2Y_1}\Delta T,
 \qquad \Delta L_2=-\Delta L_1.
\]
若第一根杆位于左边，则连接点向右的位移正是 $\Delta L_1$。连接点可以移动，所以通常并不是 $\Delta L_1=\Delta L_2=0$；一段略伸长时另一段必须等量缩短。仅当 $\alpha_1Y_1=\alpha_2Y_2$ 时，两段才都不改变长度。
分母的单位为长度除以压强，所得 $F$ 的单位为力；若两种材料相同，就回到 $F=-AY\alpha\Delta T$，与总长度无关。若仅一段非常柔软，更多的弹性压缩会集中在该段，这也解释了不能简单平均两个杨氏模量的原因。
\end{xesolution}
